% 本文件由 LaTeX 排版 Agent 根据有效 translation.json 自动生成。
% author=Apocrypha; work_id=0846; title=多默童年福音

\chapter{多默童年福音}

% structure: p0001 source=h1 target=omitted reason=page-title-already-rendered-by-parent

% structure: p0002 source=h2 target=section reason=relative-heading-rank; base=h2

\section{第一希腊语版本}

\Emph{\Person{多默}，以色列哲人，关于主童年的记述}。\par

一、我\Person{多默}，以色列人，给你们写下这篇记述，为使所有外邦中的弟兄们都知道我们的主\Person{耶稣基督}在童年所行的奇事，即祂在我们家乡出生后所行的。故事的开端如下：—\par

二、这孩子\Person{耶稣}五岁时，在山溪的渡口玩耍；祂将流动的水聚成水塘，立即使之清澈，并且仅凭话语就使水听从祂。祂又取了软泥，用它塑了十二只麻雀。祂做这些事时正是安息日。还有许多别的孩子和祂一同玩耍。有一个犹太人看见\Person{耶稣}在安息日所做的事，就立刻去告诉祂的父亲\Person{若瑟}说：\Quote{看，你的儿子在溪边，取了泥，做了十二只鸟，亵渎了安息日。}\Person{若瑟}来到现场，看见后，向祂喊道：\Quote{你为什么在安息日做不合法的事？}\Person{耶稣}拍手，向麻雀喊道：\Quote{去吧！}麻雀就飞去，叫着飞走了。犹太人看见这事，十分惊奇，就去向他们的首领报告他们所看见\Person{耶稣}做的事。\par

三、那时，经师\Person{亚纳斯}的儿子正与\Person{若瑟}站在那里；他拿了一根柳条，放掉了\Person{耶稣}所聚的水。\Person{耶稣}看见所发生的事，就很生气，对他说：\Quote{啊，邪恶、不敬神、愚昧的！水塘和水对你做了什么伤害？看，如今你要像树一样枯干，不再生出叶子、根或果实。}那男孩立刻就完全枯干了。\Person{耶稣}离开，去了\Person{若瑟}的家。但那个枯干了的男孩的父母抱起他，哀叹他的青春，把他带到\Person{若瑟}那里，并责备他说：\Quote{\Emph{他们说}，你们有这样一个孩子，做这样的事。}\par

四、此后，祂再次经过那村庄；有一个男孩跑过来撞到祂，打了祂的肩膀。\Person{耶稣}很生气，对他说：\Quote{你不得再回到你来的路上去。}他立即倒地死了。有些看见这事的人说：\Quote{这孩子是从哪里来的，以致他的话每一句都必定应验？}死去的男孩的父母到\Person{若瑟}那里，责备他说：\Quote{既然你有这样一个孩子，你就不能与我们同住在村庄里；否则就教祂祝福，而不是诅咒，因为祂正在杀死我们的孩子。}\par

五、\Person{若瑟}把小孩叫到一旁，告诫祂说：\Quote{你为什么做这样的事，以致这些人受苦，恨我们，迫害我们？}\Person{耶稣}说：\Quote{我知道这些话不是出自你；但为了你的缘故，我会保持沉默；但他们要受惩罚。}那些控告祂的人立刻瞎了。看见这事的人非常害怕，十分困惑，谈论祂说：\Quote{他所说的每一句话，无论是好是坏，都是一个行动，成了一个奇事。}他们看见\Person{耶稣}做了这样的事，\Person{若瑟}起来，抓住祂的耳朵，用力扯。那孩子非常生气，对他说：\Quote{对你来说，寻求却找不到，已经够了；而你确实没有明智地行事。难道你不知道我是属于你的吗？不要烦扰我。}\par

六、有一位老师，名叫\Person{匝凯}，站在某处，听见\Person{耶稣}这样对祂父亲说话，极其惊奇，因为一个孩子竟能如此说话。几天之后，他来见\Person{若瑟}，对他说：\Quote{你有一个聪明的孩子，他有些头脑。那么，把他交给我，让他学习字母；我将教他字母以及一切知识，包括如何称呼所有长者，尊敬他们如祖先和父亲，以及如何爱那些同龄人。}祂把从\OriginalTerm{阿尔法}{Alpha}到\OriginalTerm{奥梅加}{Omega}的所有字母都清楚地、精确地告诉了他。祂看着老师\Person{匝凯}，对他说：\Quote{对\OriginalTerm{阿尔法}{Alpha}性质一无所知的你，怎能教别人\OriginalTerm{贝塔}{Beta}？你这个伪君子！首先，如果你知道，就教A，然后我们才相信你关于B的话。}然后祂开始就第一个字母考问老师，老师不能回答祂。当着许多人的面，孩子对\Person{匝凯}说：\Quote{老师，请听第一个字母的结构，留意它如何有线条，一条中间笔划交叉那些你常见的线条；线条汇聚在一起；最高部分支撑着它们，又将它们汇聚在一个头部之下；有三个\Emph{交点}；同类的；主要的与附属的；等长的。这就是A的线条。}\par

七、当老师\Person{匝凯}听见孩子说出关于第一个字母的如此伟大比喻时，他对这样的叙述和祂的教训感到非常困惑。祂对在场的人说：\Quote{唉！我，可怜的人，很困惑，因我把这孩子带到这里而自取其辱。那么，我恳求你，\Person{若瑟}兄弟，带祂回家。我无法忍受祂目光的严厉；我完全无法理解祂的意思。那孩子不属于这个世界；祂甚至能驯服火。祂肯定是在创世之前就存在的。什么样的肚子生了祂，什么样的子宫养育了祂，我不知道。唉！我的朋友，祂把我迷住了；我抓不住祂的意思：我真是三次可怜，我欺骗了自己。我努力想要有一个学生，结果却发现我有了一个老师。我的心中充满羞愧，朋友们，因为我一个老人，竟被一个孩子打败了。除了沮丧和死亡，我没有别的待遇，因为这个孩子，因为此刻我无法正视祂；当人人都说我是被一个小孩子打败时，我能说什么呢？我如何解释祂所说的第一个字母的线条？我不知道，噢，我的朋友们；因为我既不知道祂的开始，也不知道祂的结局。因此，我恳求你，\Person{若瑟}兄弟，带祂回家。祂究竟是伟大的神还是天使，或我该说什么，我不知道。}\par

8. 当犹太人鼓励\Person{匝凯}时，孩子大声笑着说：\Quote{现在让你的学识结出果实，让心中瞎眼的人看见。我从上面来，为要诅咒他们，并召唤他们归向上面的事，正如那为我缘故差遣我的主所命令的。}孩子说完话后，所有因祂诅咒而倒下的人立刻都痊愈了。此后，没有人再敢惹祂发怒，免得祂诅咒他，使他成为残废。\par

9. 过了几天，耶稣在某一座房子的顶楼玩耍，与祂一起玩耍的一个孩子从楼上跌下来摔死了。其他孩子看见，就逃跑了，只有耶稣独自站着。死孩子的父母来了，责怪……并威胁祂。耶稣从屋顶跳下，站在孩子的尸体旁，大声喊叫说：\Quote{\Person{Zeno}——这是他的名字——起来，告诉我：是我把你推下去的吗？}他立刻站起来说：\Quote{当然不是，我的主；你没有推我下去，而是使我复活了。}看见这事的人都大为惊奇。孩子的父母因所发生的奇迹光荣天主，并朝拜了耶稣。\par

10. 过了几天，一个年轻人在角落里劈柴，斧头落下砍中他的脚底，脚掌被劈成两半，他因失血过多而死。众人大乱，人们聚拢过来，孩童耶稣也跑到那里。祂挤过人群，抓住年轻人的伤脚，他便立刻痊愈了。祂对年轻人说：\Quote{现在起来，劈柴，并记念我。}众人看见所发生的事，都朝拜孩子，说：\Quote{天主的圣神确实住在这孩子内。}\par

11. 当祂六岁时，祂母亲给了祂一个水罐，打发祂去打水带回家里。但祂在人群中撞到一个人，水罐打破了。耶稣展开祂披着的斗篷，装满水，带回给母亲。祂的母亲看见所发生的奇迹，就亲吻了祂，并把所看见祂行的奥秘藏在心里。\par

12. 又在播种时，孩子与祂的父亲出去在他们田里撒种。当祂的父亲播种时，孩童耶稣也播了一粒麦子。祂收割后，打场，得了一百个\OriginalTerm{苛尔}{kors}；祂把村子里的所有穷人都叫到打谷场，把麦子分给他们，若瑟拿走了剩下的麦子。祂行这奇迹时是八岁。\par

13. 祂的父亲是个木匠，那时做犁和轭。一个富人向他定做一张床。其中一块所谓的横木太短了，他们不知怎么办。孩童耶稣对祂的父亲若瑟说：\Quote{把那两块木头放下来，在中间对齐。}若瑟照孩子说的做了。耶稣站在另一端，抓住那块较短的木头，把它拉长，使之与另一块相等。祂的父亲若瑟看见了，就惊奇，拥抱孩子，并祝福祂说：\Quote{我有福了，因为天主赐给了我这孩子。}\par

14. 若瑟见这孩子心智和身体都很健壮，又决定不该让祂不识字，就把祂领走，交给另一位老师。老师对若瑟说：我先教祂希腊文字母，再教希伯来文。原来老师已经听说过以前对这孩子的考验，因此怕祂。他还是写下了字母表，花了很多时间注意祂，但祂没有回答他。耶稣对他说：\Quote{如果你真是老师，又精通字母，就告诉我\OriginalTerm{阿尔法}{Alpha}的力量，我便告诉你\OriginalTerm{贝塔}{Beta}的力量。}老师为此发怒，打了祂的头。孩子疼痛，就诅咒了他；他立刻昏倒，面朝下跌倒在地。孩子回到若瑟家；若瑟很悲伤，吩咐祂母亲说：\Quote{不要让祂出门，因为惹祂生气的人都得死。}\par

15. 过了一些时候，另一位老师，也是若瑟的挚友，对他说：把孩子带到我的学校来；或许我能哄祂学会识字。若瑟说：\Quote{兄弟，如果你有勇气，就带祂去吧。}他怀着恐惧和极大的痛苦把祂带走了；但孩子愉快地跟着去了。祂大胆地走进学校，发现有一本书放在书桌上；祂拿起书来，没有读上面的字母，却开口借着圣神说话，给周围的人讲授法律。一大群人聚拢来站着听，对祂教导的成熟和话语的流利感到惊奇，而且祂小小年纪，竟能如此讲话。若瑟听见这事，害怕了，跑到学校，担心这位老师也缺乏经验。老师对若瑟说：\Quote{兄弟，你知道，我收这孩子做学生，祂充满极大的恩宠和智慧；但我恳求你，兄弟，带祂回家吧。}孩子听了这话，立刻笑他，说：\Quote{因为你说了正确的话，作了正确的见证，为了你的缘故，那个被击倒的人也要痊愈。}另一位老师立刻痊愈了。若瑟便领着孩子回家去了。\par

16. 若瑟打发他的儿子雅各伯去捆柴带回家，孩童耶稣也跟随着他。雅各伯正在拾柴时，一条毒蛇咬了雅各伯的手。他疼痛难忍，快要死了，耶稣上前，向伤口吹了一口气；疼痛立刻停止，毒蛇爆裂，雅各伯转瞬安然无恙。\par

17. 此后，若瑟的一个邻居的婴儿患病死了，他的母亲痛哭不已。耶稣听见有极大的哀号与骚动，就急忙跑过去，见那孩子死了，就摸着他的胸膛说：\Quote{我对你说，孩子，不要死，要活着，并与你母亲在一起。}那孩子立刻抬眼望，笑了。耶稣又对那妇人说：\Quote{抱他，给他喂奶，并记念我。}围观的众人见了，都惊奇说：\Quote{这孩子真是天主或天主的天使，因为他每一句话都确实应验。}耶稣就从那里出去，与其他孩子玩耍。\par

18. 过了一些时候，在某处建房时发生一阵大骚动，耶稣起身来到那地方。看见一个人躺在地上死了，就拉着他的手说：\Quote{我告诉你，起来，继续你的工作。}那人立刻起来，朝拜祂。众人见了，惊奇说：\Quote{这孩子是从天上来的，因为他救了许多灵魂脱离死亡，而且他一生都继续救人。}\par

19. 当祂十二岁时，祂的父母照例与同行的人上耶路撒冷去过逾越节。逾越节后，他们又回家。在他们回家的路上，孩童耶稣转回耶路撒冷。祂的父母以为祂在同行的人群中。走了一天的路程后，他们在亲戚中找祂；找不到，就非常忧伤，转回城里寻找祂。第三天，他们在圣殿里找到祂，坐在教师们中间，一边听法律，一边向他们提问。众人都留意祂，并惊奇祂虽是个孩童，却能堵住民间长老和教师们的口，解释法律的重点和先知们的比喻。祂的母亲玛利亚上前对祂说：\Quote{孩子，你为什么这样对待我们？看，我们忧苦地找你。}耶稣对他们说：\Quote{你们为什么找我？难道你们不知道我必须在我父亲的事务上吗？}经师和法利塞人说：\Quote{你是这孩子的母亲吗？}她说：\Quote{我是。}他们对她说：\Quote{你在妇女中是有福的，因为天主祝福了你胎中的果子；因为这样的荣耀、这样的美德和智慧，我们从未见过听过。}耶稣就起来，跟随母亲，服从父母。祂的母亲把发生的这一切都默存在心。耶稣在智慧、身量和恩宠上渐长。\BibleRef{路 2:41}愿光荣归于祂，直到永远。阿们。\par

% structure: p0023 source=h2 target=section reason=relative-heading-rank; base=h2

\section{\Emph{希腊文第二式}}

\Emph{圣宗徒\Person{多默}关于主童年事迹的记述}\par

1. 我，以色列人\Person{多默}，认为有必要将我们的主耶稣基督在童年时所行的大事告知所有外邦的弟兄，那时祂居于肉身，在\Place{纳匝肋}城，年龄在第五年。\par

2. 有一天，下着暴雨，祂从母亲所在的屋里出来，在流水的土地上玩耍。祂做了水池，引水进去，水池满了水。然后祂说：\Quote{我要你们变成清澈优良的水。}它们立刻变成了。有一个男孩，是经师\Person{亚纳斯}的儿子，经过时用他拿着的柳枝把水池打垮，水流了出来。耶稣转过身来对他说：\Quote{你这不敬而邪恶的人，水池怎么得罪了你，你要倒空它们？你不得继续前行，你要像你拿着的树枝一样干枯。}他走了一会儿，不久就跌倒死了。和他一起玩耍的孩子们见了，就惊奇，跑去告诉那死孩子的父亲。他跑来，见儿子死了，就去责斥\Person{若瑟}。\par

3. 耶稣用那泥做了十二只麻雀，那天正是安息日。有一个孩子跑去告诉\Person{若瑟}说：\Quote{看，你的孩子在溪边玩耍，用泥做了麻雀，这是不合法的。}若瑟听见后，就去对那孩子说：\Quote{你为什么这样做，亵渎安息日？}耶稣却未回答他，只看着麻雀说：\Quote{去吧，飞，活着，记念我。}这话一出口，麻雀就飞起来，升到空中。\Person{若瑟}见了，就惊奇。\par

4. 几天后，当耶稣走过城中，一个男孩朝祂扔石头，打中祂的肩膀。耶稣对他说：\Quote{你不得继续前行。}他立即跌倒，也死了。恰好在场的人都大为惊讶，说：\Quote{这孩子是从哪里来的？他说的每句话都确实应验。}他们也去责斥\Person{若瑟}，说：\Quote{你不可能再与我们同住在这城里了；但如果你愿意如此，就教导你的孩子祝福，而不是诅咒；因为他杀死我们的孩子，他说的每句话都确实应验。}\par

5. 若瑟坐在自己的位上，那孩子站在他面前；他就抓住祂的耳朵，用力拧。耶稣定睛看着他，说：\Quote{对你来说已经够了。}\par

6. 次日，他拉着耶稣的手，领祂到一个名叫\Person{匝凯}的教师那里，对他说：\Quote{老师，请收下这孩子，教他认字。}他说：\Quote{兄弟，把他交给我，我来教他圣经；我要劝他祝福众人，而不是诅咒。}耶稣听见就笑了，对他们说：\Quote{你们说你们所知道的，但我比你们知道得多，因为我在万世之前。我知道你们的祖先和先祖何时出生，也知道你们寿数几何。}他们听了，大为惊讶。耶稣又对他们说：\Quote{你们惊讶，是因为我对你们说，我知道你们寿数几何。我确实知道世界何时被造。看，你们现在不信我。当你们看见我的十字架时，那时你们就会相信我说的是真理。}他们听了这些话，都大为惊讶。\par

7. 匝加乌斯用希伯来文写好字母表后，对祂说：\Quote{阿尔法。}孩子说：\Quote{阿尔法。}老师又说：\Quote{阿尔法。}孩子也同样说。然后老师第三次说阿尔法。耶稣望着老师的脸说：\Quote{你连阿尔法都不认识，怎能教别人贝塔呢？}那孩子从阿尔法开始，自己说了二十二个字母。然后又说道：\Quote{老师，请听第一个字母的顺序，知道它有多少入口、多少线条，以及共同的、交叉和汇合的笔画。}匝加乌斯听到关于一个字母这样的解释，惊讶得无法回答。他转身对若瑟说：\Quote{兄弟，这孩子肯定不属于地上。请把他从我这里带走。}\par

8. 此后有一天，耶稣与别的孩子在屋顶上玩耍。一个孩子被另一个孩子推挤，摔到地上死了。看见这事，与祂玩耍的孩子们都逃跑了，只有耶稣独自站在那孩子摔下的屋顶上。死孩子的父母听到消息，哭着跑来；发现孩子躺在地上死了，耶稣站在上面，就以为孩子是被祂推下去的；他们注视着耶稣，辱骂祂。耶稣看见这情景，立刻从屋顶下来，站在尸体头前，对祂说：\Quote{泽诺，是我把你推下去的吗？站起来，告诉我们。}这是那孩子的名字。话音一落，那孩子站起来朝拜耶稣，说：\Quote{我主，不是你推我下去的，而是你使我从死里复活了。}\par

9. 几天后，一个邻居劈柴时，斧头砍到了脚的下部，因失血而濒死。许多人跑过来，耶稣也同他们来到那里。祂抓住那年轻人的伤脚，立刻医治好，并对他说：\Quote{起来，劈你的柴吧。}他站起来朝拜祂，感谢着劈柴。同样，所有在场的人都惊奇，并向祂致谢。\par

10. 耶稣六岁时，祂的母亲玛利亚派祂去泉边取水。路上，水罐打破了。耶稣来到泉边，展开外衣，从泉中取水，装满外衣，把水带给母亲。玛利亚看见这事，惊讶不已，拥抱并亲吻了祂。\par

11. 耶稣到了八岁，一位富人吩咐若瑟为他做一张床。若瑟是个木匠。他出到田里取木料，耶稣与他同去。他砍了两块木头，用斧子刨平，将一块放在另一块旁边，一量却发现短了。他看见这事便忧愁，想再找一块。耶稣看见，对祂说：\Quote{把这两块放在一起，使两端对齐。}若瑟不明白孩子什么意思，却照做了。耶稣又对他说：\Quote{牢牢握住短的那块。}若瑟惊讶地握住了。然后耶稣也握住另一端，向自己这边拉，使它与另一块木头等长。耶稣对若瑟说：\Quote{不要再忧愁，尽管做你的活。}若瑟看见这事，非常惊奇，自言自语说：\Quote{我有福了，因为天主赐给我这样一个孩子。}他们回到城里，若瑟把这事告诉了玛利亚。她听见并看见儿子奇异的奇迹，就欢喜，光荣祂，并偕同圣父及圣神，从今世直到永远。阿们。\par

% structure: p0036 source=h2 target=section reason=relative-heading-rank; base=h2

\section{拉丁文本}

\Emph{按多玛斯所记耶稣童年史于此开始。}\par

% structure: p0038 source=h3 target=subsection reason=relative-heading-rank; base=h2

\subsection{第一章 玛利亚和若瑟如何带祂逃往埃及。}

当黑落德搜寻我们主耶稣基督要杀害祂，引起骚动时，一位天使对若瑟说：\Quote{带着玛利亚和她的孩子，逃往埃及，躲避那些要杀祂的人。}耶稣进入埃及时两岁。\par

当祂走过一片麦田时，祂伸出手，摘下麦穗，放在火上烤，搓着，开始吃。\par

他们进入埃及后，在一寡妇家得到款待，在那里住了一年。\par

耶稣三岁时，看见孩子们玩耍，就与他们一起玩。祂拿起一条干鱼，放进盆里，命令它游动。鱼就游动起来。祂又对鱼说：\Quote{倒出你里面的盐，走进水里。}事情就这样成就了。邻居们看见所发生的事，就告诉了那寡妇，玛利亚——祂的母亲——住在她家。寡妇一听，就急忙将她们赶出了家门。\par

% structure: p0043 source=h3 target=subsection reason=relative-heading-rank; base=h2

\subsection{第二章 一位老师如何将祂赶出城。}

当耶稣与祂的母亲玛利亚穿过城中市场时，祂看见一位老师正在教学生。看啊，十二只麻雀吵闹着，从墙上掉进那正在教书的老师怀里。耶稣看见，十分开心，就站住了。那老师看见祂嬉笑，就怒气冲冲地对学生说：\Quote{去把他给我带来。}他们带祂到老师面前，老师揪住祂的耳朵说：\Quote{你看见了什么，让你这么开心？}耶稣对他说：\Quote{老师，看，我手里满是小麦。我拿给他们看，把小麦撒在他们中间，他们就把小麦从危险的街上运出去；为此他们互相争斗分小麦。}耶稣没有离开那地方，直到事情成就。之后，那老师开始把祂和祂的母亲赶出城。\par

% structure: p0045 source=h3 target=subsection reason=relative-heading-rank; base=h2

\subsection{第三章 耶稣如何离开埃及。}

看，主的天使遇见\Person{玛利亚}，对她说：\Quote{抱起这孩子，回到犹太人的地方去，因为那些寻索他性命的人已经死了。}\Person{玛利亚}便带着\Person{耶稣}起身，前往她父亲产业中的\Place{纳匝肋}城。当\Person{黑落德}死后，\Person{若瑟}离开\Place{埃及}，将\Person{耶稣}保存在旷野中，直到\Place{耶路撒冷}那些寻索这孩子性命的人安静下来。他感谢天主，因为天主赐给了他智慧，也因他在上主天主面前蒙了恩宠。阿们。\par

% structure: p0047 source=h3 target=subsection reason=relative-heading-rank; base=h2

\subsection{第四章 主耶稣在\Place{纳匝肋}城所行的事。}

身为以色列人和主的宗徒的\Person{多默}，荣耀地记述了\Person{耶稣}从\Place{埃及}出来进入\Place{纳匝肋}后所行的作为。我至亲爱的弟兄们，请你们都明白主耶稣在\Place{纳匝肋}城时所行的事，其第一章如下：——\par

当\Person{耶稣}五岁时，地上降下大雨，孩童\Person{耶稣}在雨中走来走去。雨势凶猛，祂将雨水收集到一个鱼池中，并用自己的话命令水变得清澈。水立刻变得清澈。祂又取那鱼池中的泥土，做了十二只麻雀。\Person{耶稣}在犹太人的孩童中做这事时，正是安息日。犹太人的孩童便去告诉祂的父亲\Person{若瑟}说：\Quote{看，你的儿子和我们一起玩耍，他拿泥做了麻雀，这在安息日是不允许的；他犯了安息日。}\Person{若瑟}来到孩童\Person{耶稣}面前，对祂说：\Quote{你为什么做了在安息日不允许做的事？}\Person{耶稣}张开双手，命令麻雀说：\Quote{飞到空中去吧，无人会杀你们。}麻雀便飞起来，开始鸣叫，赞美全能的天主。犹太人看见所发生的事，十分惊奇，便去传扬\Person{耶稣}所行的奇迹。但有一个与\Person{耶稣}同处的法利塞人，拿起一根橄榄枝，开始放掉\Person{耶稣}所造的水泉里的水。\Person{耶稣}看见这事，便愤怒地对他说：\Quote{你这不虔敬又无知的索多玛人，我的工程——水泉——对你做了什么伤害？看，你要变成一棵枯树，没有根，没有叶，也没有果实。}那人立刻枯干，倒在地上死了。他的父母把尸体抬走，责备\Person{若瑟}说：\Quote{看你儿子做了什么；教导他祈祷，不要亵渎。}\par

% structure: p0050 source=h3 target=subsection reason=relative-heading-rank; base=h2

\subsection{第五章 市民因耶稣的所作所为向\Person{若瑟}发怒。}

几天后，\Person{耶稣}与\Person{若瑟}在城中行走时，有一个孩子跑上来击打\Person{耶稣}的胳膊。\Person{耶稣}对他说：\Quote{你走不完你的路程了。}孩子立刻倒在地上死了。看见这些奇事的人喊道：\Quote{这孩子是从哪里来的？}他们对\Person{若瑟}说：\Quote{这样的孩子不该在我们中间。}\Person{若瑟}去把祂带来。他们对他说：\Quote{离开这个地方；但如果你必须和我们同住，就教导他祈祷，不要亵渎；因为我们的孩子已被杀了。}\Person{若瑟}叫来\Person{耶稣}，责备祂说：\Quote{你为什么亵渎？因为住在这里的人恨我们。}\Person{耶稣}说：\Quote{我知道这些话不是我的，而是你的；但为了你，我会保持沉默；让他们凭智慧去省察吧。}那些反对\Person{耶稣}的人立刻瞎了眼。他们走来走去，说：\Quote{从他口中所出的一切话都成就了。}\Person{若瑟}看见\Person{耶稣}所做的事，愤怒地抓住祂的耳朵；\Person{耶稣}愤怒地对\Person{若瑟}说：\Quote{你看见我就够了，不要碰我。因为你不知道我是谁；但如果你知道，就不会惹我发怒。即使现在我和你们在一起，我却在你们之前就存在了。}\par

% structure: p0052 source=h3 target=subsection reason=relative-heading-rank; base=h2

\subsection{第六章 耶稣如何被老师对待。}

有一个名叫\Person{匝凯}的人，听了\Person{耶稣}对\Person{若瑟}所说的一切话，十分惊讶，自言自语说：\Quote{从未见过这样说话的孩子。}他走到\Person{若瑟}面前说：\Quote{你有一个聪明的孩子；把他交给我，让他学习字母；当他完全学会了字母，我就体面地教导他，使他不再愚蠢。}\Person{若瑟}回答说：\Quote{除了天主，没有人能教导他。你不相信那个小孩会是无足轻重的吗？}\Person{耶稣}听见\Person{若瑟}这样说，就对\Person{匝凯}说：\Quote{夫子，凡从我口出来的都是真实的。在万有之先，我就是主，而你们是外邦人。万世的荣耀已赐给了我，你们却什么也没有得到；因为我在万世之先就已存在。我知道你将有多少年岁的寿命，且你将被掳到异地：这是我的圣父所指定的，为叫你们明白，凡从我口出来的都是真实的。}站在旁边的犹太人听见\Person{耶稣}所说的话，十分惊奇，说：\Quote{我们看见了这样奇妙的事，从那个孩子听见了这样的话语，是我们从未听见过，也从未从任何其他人——无论是大司祭、经师还是法利塞人——那里可能听见的。}\Person{耶稣}回答说：\Quote{你们为什么惊奇？你们认为我说了实话是难以置信的吗？我知道你们和你们的祖先何时出生，更告诉你们，当世界被造时，我也知道是谁派遣我到你们这里来。}犹太人听见这孩子所说的话，都惊奇，因为他们不能回答。那孩子自言自语，欢跃说：\Quote{我对你们说了个比喻；我知道你们是软弱无知的。}\par

那位老师对若瑟说：\Quote{把他带到这里来，我要教他识字。}若瑟就牵着孩童耶稣，领他到一所老师家中，那里有其他孩子正在学习。老师用温和的话语开始教祂字母，并为祂写下第一行，即从A到T，然后开始拍打祂、教导祂。那位老师打了孩子的头；孩子挨打后对他说：\Quote{我应当教你，而不是你教我；你所想教我的字母我都知道，而且我知道，你对我来说就像器皿，只发出声音，却没有智慧。}接着，祂开始背诵那一行，完整而快速地念出从A到T的字母。祂看着老师，对他说：\Quote{你连A和B是什么都无法告诉我们，又怎能想教导别人？伪君子啊，你若知道并告诉我A，我就告诉你B。}当那位老师开始讲述第一个字母时，他无法给出任何回答。耶稣对匝凯说：\Quote{老师，请听我说：明白第一个字母。看它如何有两条线；向中间推进、静止、给予、散开、变化、威胁；三重与双重交织；同时同质，一切共有。}\par

匝凯见祂如此分解第一个字母，对这第一个字母、对这样的人和这样的学识感到惊愕；他喊道：\Quote{祸哉，我完全惊愕了；我因这孩子自取其辱。}他对若瑟说：\Quote{我恳求你，兄弟，把他从我这里带走，因为我不能看他的脸，也不能听他的大能话语。因为这孩子能制服烈火、勒住大海：他是在万世之前诞生的。什么子宫怀了他？什么母亲哺育了他？我不知道。哦，我的朋友们，我神志不清；我成了可怜的笑柄。我曾以为我收了一个学生，但他却成了我的老师。我的羞辱无法摆脱，因为我老了；我找不到话对他说。我所能做的只有患上重病，离开这个世界，或者离开这城，因为众人都见到了我的羞辱。一个婴孩欺骗了我。我能给别人什么回答？或能说什么话？因为他在第一个字母上胜过了我。我哑口无言，哦，我的朋友们和相识们；我找不到回答他的开头或结尾。现在我恳求你，若瑟兄弟，把他从我这里带走，领他回家，因为他是老师，或是主，或是天使。我不知道说什么。}耶稣转向与匝凯在一起的犹太人，对他们说：\Quote{让看不见的看见，让不理解的明白；让聋子听见，让因我而死的人复活；让那些高升的，我召他们到更高之处，正如那派我到你处者所吩咐的。}耶稣说完后，所有因祂的话而患任何疾病的人都痊愈了。他们不敢对他说话。\par

% structure: p0056 source=h3 target=subsection reason=relative-heading-rank; base=h2

\subsection{第七章 耶稣如何使一个男孩复活}

有一天，耶稣与孩子们在一所房子上玩耍，祂开始与他们嬉戏。其中一个男孩从后门掉下去，立即死了。孩子们看见后都逃跑了，但耶稣留在房子里。死者的父母来到后，指责耶稣说：\Quote{一定是你使他掉下去的。}他们辱骂祂。耶稣从房子上下来，站在死去的孩子面前，大声呼喊那孩子的名字：\Quote{Sinoo，Sinoo，起来说，是不是我使你掉下去的。}他立刻起来说：\Quote{不是，我的主。}他的父母看见耶稣行了如此大的奇迹，就光荣天主，并朝拜耶稣。\par

% structure: p0058 source=h3 target=subsection reason=relative-heading-rank; base=h2

\subsection{第八章 耶稣如何治愈一个男孩的脚}

几天后，那城中的一个男孩劈柴时伤了自己的脚。一大群人走去他那里，耶稣也和他们一起去。祂触摸受伤的脚，脚立刻痊愈了。耶稣对他说：\Quote{起来，劈柴，并记念我。}众人看见祂所行的奇迹，就朝拜耶稣，说：\Quote{我们确实确信你是天主。}\par

% structure: p0060 source=h3 target=subsection reason=relative-heading-rank; base=h2

\subsection{第九章 耶稣如何用外衣取水}

耶稣六岁时，祂的母亲派祂去打水。耶稣来到泉边（或井边），那里有大批人群，他们打碎了祂的水罐。祂脱下所穿的外衣，盛满水，带回给母亲玛利亚。祂的母亲看见耶稣所行的奇迹，亲吻祂说：\Quote{主啊，求你听我，拯救我的儿子。}\par

% structure: p0062 source=h3 target=subsection reason=relative-heading-rank; base=h2

\subsection{第十章 耶稣如何播种小麦}

播种时节，若瑟出去播种小麦，耶稣跟随着他。若瑟开始播种时，耶稣伸出手，抓了一把小麦，撒了出去。收割时，若瑟来收庄稼。耶稣也来了，收集祂所撒的麦穗，他们获得了一百斗的上好谷粒；祂叫来穷人、寡妇和孤儿，将祂所收获的小麦分给他们。若瑟也拿了少许同批小麦，作为耶稣对家室的祝福。\par

% structure: p0064 source=h3 target=subsection reason=relative-heading-rank; base=h2

\subsection{第十一章 耶稣如何使短木料与长木料一样长}

耶稣长到八岁。若瑟是位工匠师傅，惯常制作犁和牛轭。有一天，一个富人对若瑟说：\Quote{师傅，请给我做一张床，既要实用又要美观。}若瑟很苦恼，因为他拿来作工的木料太短了。耶稣对他说：\Quote{不要烦恼。你握住这木料的一端，我握住另一端，让我们把它拉长。}他们照做了，木料立刻适合了他的意图。祂对若瑟说：\Quote{看，做你想做的工作吧。}若瑟看见所成的事，拥抱祂说：\Quote{我有福了，因为天主赐给我这样的儿子。}\par

% structure: p0066 source=h3 target=subsection reason=relative-heading-rank; base=h2

\subsection{第十二章 耶稣如何被交给老师学习字母}

若瑟见祂如此蒙恩宠，身材日益增长，便认为应当送祂去学习字母。于是将祂交给另一位老师教导。那老师对若瑟说：\Quote{你希望我教那孩子什么字母？}若瑟回答说：\Quote{先教祂外邦字母，再教希伯来字母。}那老师知道祂极聪慧，便欣然收下祂。他为祂写下第一行，即 A 和 B，教了几个时辰。但耶稣默不作声，不回答他。耶稣对老师说：\Quote{你若真是一位老师，若你真认识字母，就告诉我 A 的能力，我便告诉你 B 的能力。}那老师于是勃然大怒，击打祂的头。耶稣发怒，诅咒了他；他立刻跌倒，死了。\par

耶稣回到家中。若瑟吩咐祂的母亲玛利亚，不可让祂离开自家的庭院。\par

% structure: p0069 source=h3 target=subsection reason=relative-heading-rank; base=h2

\subsection{第十三章　祂如何被交给另一位老师}

许多天后，又来了一位老师，是若瑟的朋友，对若瑟说：\Quote{把祂交给我，我必用极大的温柔教祂字母。}若瑟对他说：\Quote{你若能，就带祂去教吧。愿此事伴随着喜乐。}那老师领着祂，心中畏惧而极其坚定，欣然牵着祂。及至来到老师家，祂发现那里放着一本书，便拿起书打开，却没有读书中所写的；祂却开口，藉着圣神发言，教导法律。当时所有站着的人都专心地听祂；老师坐在祂旁边，愉快地听着，并恳请祂多教导他们。又有一大群人聚集，听到祂所教导的全部神圣教训，以及从祂口中说出的优美言辞——祂虽为孩童，却说出这样的事。若瑟听见这事，便害怕，跑去找那老师，老师正与耶稣在一起，对若瑟说：\Quote{兄弟，你知道我已接受你的孩子，要教导或训练祂；但祂充满极大的庄重和智慧。看，现在你带祂回家，欢喜吧，我的兄弟；因为祂所拥有的庄重是主所赐的。}耶稣听见老师这样说，便高兴起来，说：\Quote{看，现在老师，你说了真话。为了你，那死去的人将要复活。}若瑟便带祂回家。\par

% structure: p0071 source=h3 target=subsection reason=relative-heading-rank; base=h2

\subsection{第十四章　耶稣如何从蛇咬中解救雅各伯}

若瑟差遣雅各伯去拾柴，耶稣跟随着他。雅各伯正拾柴时，一条毒蛇咬了他；他倒在地上，如同中毒而死。耶稣看见，便向他的伤口吹气；雅各伯立刻痊愈，毒蛇死了。\par

% structure: p0073 source=h3 target=subsection reason=relative-heading-rank; base=h2

\subsection{第十五章　耶稣如何使一个孩子复活}

数日后，一个邻家的孩子死了，他的母亲极为哀痛。耶稣听见，便走去站在孩子上面，拍打他的胸膛，说：\Quote{我对你说，孩子，不要死，要活过来。}孩子立刻起来了。耶稣对孩子的母亲说：\Quote{接过你的儿子，给他喂奶，并记念我。}群众看见这奇迹，说：\Quote{这孩子确是从天来的；因为祂已使许多灵魂脱离死亡，并使所有盼望祂的人都得痊愈。}\par

经师和法利塞人对玛利亚说：\Quote{你是这孩子的母亲吗？}玛利亚说：\Quote{我确实是。}他们对她说：\Quote{你在妇女中当受赞美，\BibleRef{路 1:28}因为天主赐福了你胎中的果实，赐给你如此荣耀的孩子和如此智慧的恩赐，是我们从未见过或听过的。}耶稣起身跟随祂的母亲。玛利亚将耶稣在民间所行的一切伟大奇迹，就是治好许多患病的人，都默存在心中。耶稣在智慧和身量上渐长；凡看见祂的都光荣全能的天主圣父，祂是永受赞美的，直到永远。阿们。\par

凡此种种，我以色列人多默所见的，我都写下来，并向外邦人和我们的弟兄们讲述，以及耶稣在犹大境内所行的许多其他事。看，以色列全家都看见了这一切，从起初直到末后；耶稣在他们中间行了何等伟大的征兆和奇迹，都是极其美好的，而他们的父辈却视而不见，正如圣经所记载的，众先知也为祂在以色列万民中的作为作证。祂要按照永生旨意审判世界，因为祂是天主子，遍及普世。愿一切光荣和尊崇归于祂，直到永远，祂是永生天主，统治万世万代。阿们。\par

\SourceRecord{Translated by Alexander Walker.；Ante-Nicene Fathers；Vol. 8.；Edited by Alexander Roberts, James Donaldson, and A. Cleveland Coxe.；Buffalo, NY: Christian Literature Publishing Co.,；1886.；<http://www.newadvent.org/fathers/0846.htm>.}
